\subsection{Extension of morphisms of complexes}

Lemma 1. --- Consider a diagram of A-modules and homomorphisms

\[
\begin{array}{c} \mathbf {M} ^ {\prime} \xrightarrow {\alpha^ {\prime}} \mathbf {M} \xrightarrow {\alpha} \mathbf {M} ^ {\prime \prime} \\ f ^ {\prime} \Biggl \downarrow \quad \swarrow k / f \Biggl \downarrow \quad \swarrow k ^ {\prime \prime} \\ \mathbf {N} ^ {\prime} \xrightarrow [ \beta^ {\prime} ]{} \mathbf {N} \xrightarrow [ \beta ]{} \mathbf {N} ^ {\prime \prime} \end{array}
\]

such that \(f \circ \alpha' = \beta \circ f', \alpha \circ \alpha' = 0\) , Ker \(\beta = \operatorname{Im} \beta'\) , and \(f = k'' \circ \alpha + \beta \circ k\) and where \(\mathbf{M}'\) is projective. There exists an A-homomorphism \(k': \mathbf{M}' \to \mathbf{N}'\) such that \(f' = k \circ \alpha' + \beta' \circ k'\) .

Indeed, set \(g = f' - k \circ \alpha'\) ; we have

\[
\beta \circ g = \beta \circ f ^ {\prime} - \beta \circ k \circ \alpha^ {\prime} = f \circ \alpha^ {\prime} - \beta \circ k \circ \alpha^ {\prime} = k ^ {\prime \prime} \circ \alpha \circ \alpha^ {\prime} = 0.
\]

This implies \(\operatorname{Im}(g) \subset \operatorname{Ker}(\beta) = \operatorname{Im}(\beta')\) . Since \(\mathbf{M}'\) is projective, there therefore exists an A-homomorphism \(k': \mathbf{M}' \to \mathbf{N}'\) such that \(\beta' \circ k' = g\) , whence the lemma.

Lemma 2. --- If in the commutative diagram of A-modules and homomorphisms

\[
\begin{array}{c c c c c} \mathbf {M} ^ {\prime} & \xrightarrow {\alpha^ {\prime}} & \mathbf {M} & \xrightarrow {\alpha} & \mathbf {M} ^ {\prime \prime} \\ & & \Big \downarrow^ {u} & & \Big \downarrow^ {u ^ {\prime \prime}} \\ \mathbf {N} ^ {\prime} & \xrightarrow {\beta^ {\prime}} & \mathbf {N} & \xrightarrow {\beta} & \mathbf {N} ^ {\prime \prime} \end{array}
\]

we have \(\alpha\circ\alpha'\) =0, Ker \(\beta=\operatorname{Im}\beta'\) and if \(\mathbf{M}'\) is projective, there exists an A-homomorphism \(u':\mathbf{M}'\to\mathbf{N}'\) such that \(\beta'\circ u'=u\circ\alpha'\) .

It suffices to set \(k'' = u'', k = -u, f = 0, f' = 0\) and \(u' = k'\) in Lemma 1.

Lemma 1 bis. --- Consider a diagram of A-modules and homomorphisms

\[
\begin{array}{c} \mathbf {M} ^ {\prime} \xrightarrow {\alpha^ {\prime}} \mathbf {M} \xrightarrow {\alpha} \mathbf {M} ^ {\prime \prime} \\ \Big \downarrow_ {k ^ {\prime}} \Big \downarrow_ {f ^ {\prime}} \Big \downarrow_ {k} \Big \downarrow_ {f} \\ \mathbf {N} ^ {\prime} \xrightarrow [ \beta^ {\prime} ]{} \mathbf {N} \xrightarrow [ \beta ]{} \mathbf {N} ^ {\prime \prime} \end{array}
\]

such that \(f \circ \alpha' = \beta \circ f'\) , Ker \(\alpha = \operatorname{Im} \alpha'\) , \(\beta \circ \beta' = 0\) , and \(f' = k \circ \alpha' + \beta' \circ k'\) and where \(N''\) is injective. There exists an A-homomorphism \(k'' : M'' \to N''\) such that \(f = k'' \circ \alpha + \beta \circ k\) . Indeed, set \(q = f - \beta \circ k\) , we have

\[
g \circ \alpha^ {\prime} = f \circ \alpha^ {\prime} - \beta \circ k \circ \alpha^ {\prime} = \beta \circ f ^ {\prime} - \beta \circ k \circ \alpha^ {\prime} = \beta \circ \beta^ {\prime} \circ k ^ {\prime} = 0.
\]

This implies Ker \(g \supset \operatorname{Im} \alpha' = \operatorname{Ker} \alpha\) . Since \(\mathbf{N}''\) is injective, there therefore exists (X, p.~16, remark) an A-homomorphism \(k': \mathbf{M}'' \to \mathbf{N}''\) such that \(g = k'' \circ \alpha\) , whence the lemma.

Lemma 2 bis. --- If, in the commutative diagram of A-modules and homomorphisms

\[
\begin{array}{c c c c c} \mathbf {M} ^ {\prime} & \stackrel {{\alpha^ {\prime}}} {{\longrightarrow}} & \mathbf {M} & \stackrel {{\alpha}} {{\longrightarrow}} & \mathbf {M} ^ {\prime \prime} \\ u ^ {\prime} \Big \downarrow & & u \Big \downarrow & & \\ \mathbf {N} ^ {\prime} & \stackrel {{\beta^ {\prime}}} {{\longrightarrow}} & \mathbf {N} & \stackrel {{\beta}} {{\longrightarrow}} & \mathbf {N} ^ {\prime \prime}, \end{array}
\]

we have \(\operatorname{Ker} \alpha = \operatorname{Im} \alpha'\) , \(\beta \circ \beta' = 0\) and if \(\mathbf{N}''\) is injective, there exists an A-homomorphism \(u'' : \mathbf{M}'' \to \mathbf{N}''\) such that \(u'' \circ \alpha = \beta \circ u\) .

It suffices to set \(u' = k', u = -k, f = 0, f' = 0\) and \(k'' = u''\) in Lemma 1 bis.

PROPOSITION 1. --- Let (P, \(d_{\mathbb{P}}\) ) and (E, \(d_{\mathbb{E}}\) ) be two complexes of A-modules and r an integer.

\begin{enumerate}
\def\labelenumi{\alph{enumi})}
\item
  Let \((u_i: \mathrm{P}_i \to \mathrm{E}_i)_{i \leqslant r}\) be a family of homomorphisms such that \(d_{\mathrm{E}} \circ u_i = u_{i-1} \circ d_P\) for \(i \leqslant r\) . Suppose that \(\mathrm{P}_i\) is projective for \(i > r\) and that \(\mathrm{H}_i(\mathrm{E}) = 0\) for \(i \geqslant r\) . Then the family of the \(u_i\) extends to a morphism of complexes from P to E; two such extensions are homotopic.
\item
  Let \((u^{i}:\mathbf{P}^{i}\to\mathbf{E}^{i})_{i\leqslant r}\) be a family of homomorphisms such that \(u^{i}\circ d_{\mathrm{P}}=d_{\mathrm{E}}\circ u^{i-1}\) for \(i\leqslant r\) . Suppose that \(\mathbf{E}^{i}\) is injective for \(i>r\) and that \(\mathbf{H}^{i}(\mathbf{P})=0\) for \(i\geqslant r\) . Then the family of the \(u^{i}\) extends to a morphism of complexes from \(\mathbf{P}\) to \(\mathbf{E}\) ; two such extensions are homotopic.
\end{enumerate}

Let us prove a). The existence of an extension v of the family \((u_{i})_{i\leqslant r}\) follows immediately from Lemma 2 by induction. Let \(v'\) be another extension; set \(f=v'-v\) , and construct by induction on the integer n a homomorphism \(k_{n}:P_{n}\to E_{n+1}\) such that \(f_{n}=d_{E}\circ k_{n}+k_{n-1}\circ d_{P}\) . For \(i\leqslant r\) , take \(k_{i}=0\) . Let \(n\geqslant r\) and suppose the \(k_{i}\) constructed for \(i\leqslant n\) . Consider then the diagram

\[
\begin{array}{c} \mathrm{P} _ {n + 1} \xrightarrow {d _ {\mathrm{p}}} \mathrm{P} _ {n} \xrightarrow {d _ {\mathrm{p}}} \mathrm{P} _ {n - 1} \\ f _ {n + 1} \Biggl \downarrow \quad k _ {n} \Biggl \downarrow f _ {n} \Biggl \downarrow k _ {n - 1} \\ \mathrm{E} _ {n + 2} \xrightarrow {d _ {\mathrm{E}}} \mathrm{E} _ {n + 1} \xrightarrow {d _ {\mathrm{E}}} \mathrm{E} _ {n}. \end{array}
\]

The hypotheses of Lemma 1 are satisfied; there therefore exists an A-homomorphism \(k_{n+1}: \mathbf{P}_{n+1} \to \mathbf{E}_{n+2}\) such that \(f_{n+1} = d_{\mathbf{E}} \circ k_{n+1} + k_n \circ d_p\) , whence \(a\) ).

The proof of b) is analogous, via Lemmas 1 bis and 2 bis.

